\documentclass[11pt]{article}

\usepackage[margin=2.4cm]{geometry}
\usepackage[utf8]{inputenc}
\usepackage[T1]{fontenc}
\usepackage{lmodern}
\usepackage{amsmath}
\usepackage{amssymb}
\usepackage{booktabs}
\usepackage{array}
\usepackage{graphicx}
\usepackage{adjustbox}
\usepackage{xcolor}
\usepackage{csquotes}
\usepackage{pgfplots}
\pgfplotsset{compat=1.18}
\usepackage[colorlinks=true,linkcolor=blue,citecolor=blue,urlcolor=blue]{hyperref}
\usepackage[capitalise]{cleveref}
\usepackage{url}
\usepackage{float}
\usepackage{enumitem}

\definecolor{darkgreen}{rgb}{0.0, 0.5, 0.0}

% Same visual conventions as the manuscript: one shared axis style, legend
% outside the axis so it never covers a curve, and one colour/marker/dash per
% algorithm, identical in every figure.
\pgfplotsset{pcfplot/.style={
  tick label style={font=\scriptsize}, label style={font=\small},
  legend style={font=\scriptsize, at={(0.5,1.03)}, anchor=south,
                legend columns=2, draw=none, fill=none,
                /tikz/every even column/.append style={column sep=8pt}},
  grid=major, grid style={gray!25}}}
\pgfplotsset{
  pac/.style   ={color=red!80,          mark=square*,   thick},
  dpac/.style  ={color=red!80,          mark=diamond*,  thick, dashed},
  hpac/.style  ={color=darkgreen,       mark=*,         thick},
  dhpac/.style ={color=darkgreen,       mark=pentagon*, thick, dashed},
  rac/.style   ={color=blue!70,         mark=triangle*, thick},
  hipac/.style ={color=orange!90!black, mark=triangle*, thick},
  hopac/.style ={color=violet,          mark=square*,   thick}}

\newcommand{\tab}[1]{\input{../results/tables/#1.tex}}
\newcommand{\dat}[1]{../results/tables/rep_plot_#1.dat}

\title{Computational Study of Algorithms for the Parametric\\
Maximum Closure Problem on Directed Forests\\[0.3em]
\large Technical Report}
\author{Valerio Dose \and Fabio Furini \and Marco Locatelli}
\date{September 2026}

\begin{document}
\maketitle

\begin{abstract}
\noindent
This is a companion technical report to the manuscript \emph{``On parametric
Maximum Closure Problems over precedence forests''} (Dose, Furini and
Locatelli). The manuscript introduces the algorithms studied here
(\texttt{PaC}, \texttt{DPaC}, \texttt{HPaC}, \texttt{DHPaC}, \texttt{HIPaC},
\texttt{HOPaC} and \texttt{RaC}) with their correctness proofs and complexity
analysis, and reports a compact selection of the experiments. This report
contains the full computational study: instance design, validation
methodology, measurement protocol, campaign definitions, and complete
results broken down \emph{by size, by density and by coefficient family} --
together with the implementation analysis and the paired-ratio dispersion
that the manuscript has no room for. Every number below is generated from the
raw benchmark data of release \texttt{v0.3.0} by the repository's own
pipeline; none is hand-typed.
\end{abstract}

\section{Scope and relationship to the manuscript}

The manuscript is the reference for the problem, the algorithms and their
analysis. This report is the reference for \emph{how} the experiments were run
and \emph{what exactly} was measured, and it reports the breakdowns that the
manuscript summarizes in one sentence or omits. \Cref{sec:recap} recaps
problem and algorithms; \cref{sec:instances,sec:validation,sec:protocol,%
sec:campaigns} document instances, correctness, protocol and campaign design;
\cref{sec:results} reports every result.

The whole study was rerun from scratch in a single sweep on 2026-08-31/09-01
from one frozen commit, under the preregistered plan
\texttt{docs/EXPERIMENTAL\_PLAN\_V3.md}. All earlier benchmark results are
superseded and are cited nowhere in this report.

\section{Problem and algorithms recap}\label{sec:recap}

Let $G=(V,A)$ be a directed acyclic graph on $n$ vertices, where arc
$(i,j)\in A$ means that including vertex $i$ forces the inclusion of vertex
$j$. Each vertex has an integer profit $p_i$ (possibly negative) and a
strictly positive integer weight $w_i$. For a closed subset $S$, let
$P(S)=\sum_{i\in S}p_i$ and $W(S)=\sum_{i\in S}w_i$. The parametric Maximum
Closure Problem asks for the set of optimal solutions of
$u(\lambda):=\max_{S\text{ closed}}P(S)-\lambda W(S)$ as a function of
$\lambda$. Function $u$ is convex piecewise linear with at most $n$
breakpoints; the optimal solutions at consecutive breakpoints are nested,
which partitions $V$ into $k$ \emph{closure layers}
$\mathcal I_1,\dots,\mathcal I_k$, strictly decreasing in profit-to-weight
ratio. Computing that ordered partition is the task studied here, on directed
forests (underlying undirected graph a forest, orientations arbitrary).

\begin{description}[leftmargin=2.1cm,style=nextline]
  \item[\texttt{PaC}] Peel-and-Contract, direct-scan reference algorithm:
  $\mathcal O(n^2)$ time, $\mathcal O(n)$ space. Repeatedly peels the current
  best final vertices or contracts an arc, selecting the maximum-ratio
  candidate by a linear scan.
  \item[\texttt{DPaC}] dual of \texttt{PaC}: layers in increasing ratio order,
  same complexity.
  \item[\texttt{HPaC}] heap-based variant: the scans are replaced by two
  candidate heaps. $\mathcal O(n^2\log n)$ worst case, $\mathcal O(n)$ space.
  \item[\texttt{DHPaC}] dual of \texttt{HPaC}.
  \item[\texttt{HIPaC}/\texttt{HOPaC}] specialized heap variants for forests
  of in-trees (out-degree $\le 1$) or out-trees (in-degree $\le 1$):
  $\mathcal O(n\log n)$ time.
  \item[\texttt{RaC}] Rake-and-Compress, top-tree algorithm:
  $\mathcal O(n\log n)$ time and space, any orientation.
  \item[\texttt{BPPF}] Bounded-Precision Parametric Pseudoflow, unmodified
  upstream implementation for general precedence graphs; external baseline
  (\cref{sec:bppf}).
\end{description}

All algorithms of this repository return the same canonical object -- the
ordered sequence of closure layers with exact rational thresholds -- computed
with exact integer arithmetic: 64-bit integer coefficients, 128-bit
cross-multiplications for every ratio comparison, no floating point in any
decision path.

\section{Instance generation}\label{sec:instances}

Instances are stored in a capacity-free text format (\texttt{.pcf}): the arc
list of a forest on $n$ vertices, one integer profit and one strictly positive
integer weight per vertex, arc $(u,v)$ encoding $x_u\le x_v$. Every instance
is produced by a seeded deterministic generator -- never hand-edited -- and is
determined by a topology, a coefficient family, a size and a seed. Each is
manifested with its SHA-256 checksum, topology classification and coefficient
bounds (\texttt{instances/manifests/}), so any regenerated or downloaded
archive is verifiable byte-for-byte.

\paragraph{Affine-coefficient families.} Six families, chosen for the shape of
the ratios $p_i/w_i$ and never after an external classification.
\texttt{independent-positive}: $w_i,p_i$ independent uniform on
$\{1,\dots,1000\}$. \texttt{independent-signed}: as above with $p_i$ uniform
on $\{-1000,\dots,1000\}$. \texttt{correlated}: $p_i=w_i+\delta_i$,
$\delta_i$ uniform on $\{-50,\dots,50\}$. \texttt{anti-correlated}:
$p_i=(1001-w_i)+\delta_i$. \texttt{near-ties}: one target ratio
$\mathrm{num}/\mathrm{den}$ per instance ($\mathrm{num},\mathrm{den}$ uniform
on $\{1,\dots,20\}$), then $t_i$ uniform on $\{1,\dots,50\}$,
$w_i=\mathrm{den}\cdot t_i$, $p_i=\mathrm{num}\cdot t_i+\jmath_i$ with jitter
$\jmath_i$ uniform on $\{-2,-1,1,2\}$. \texttt{exact-ties}:
$g=\max(1,\lfloor n/20\rfloor)$ groups with their own exact ratio, vertex $i$
assigned to group $i\bmod g$, so vertices sharing a threshold are spread over
the graph rather than adjacent.

The last two families are the ones that stress exact arithmetic: they place
thresholds at infinitesimal distance (\texttt{near-ties}) or exactly equal
(\texttt{exact-ties}). They also produce markedly fewer closure layers, which
is the single most useful fact for reading absolute times, since the work of
every algorithm scales with the number of breakpoints
(\cref{tab:layers}): at $n\ge 10\,000$ the median count ranges from
$1\,222$ (\texttt{near-ties}) to $33\,183$ (\texttt{independent-signed}), a
factor of 27 between families of the same size.

\begin{table}[t]
\centering
\caption{Median number of closure layers per coefficient family, i.e.\ how
much parametric structure each family creates. This drives the running time
of every algorithm and, in \cref{sec:bppf}, the number of probes handed to
\texttt{BPPF}.}
\label{tab:layers}
\small
\tab{rep_layers_by_family}
\end{table}

\paragraph{Topologies.} \texttt{mixed-forest}: each vertex
$v\in\{2,\dots,n\}$ is connected to a uniformly chosen predecessor $u<v$
independently with probability $\varrho$, and each edge is oriented $(u,v)$ or
$(v,u)$ with probability $1/2$. With $\varrho\in\{0.3,0.6,0.9,1.0\}$ the
forest has $\varrho(n-1)$ arcs and $n-\varrho(n-1)$ components in expectation,
so $\varrho=1.0$ gives a spanning tree and $\varrho=0.3$ a forest of many
small trees. \texttt{in-forest}: vertices scanned $u\in\{1,\dots,n-1\}$, with
probability $\varrho$ one arc $(u,j)$ with $j$ uniform on $\{u{+}1,\dots,n\}$,
giving out-degree at most one; \texttt{out-forest} symmetric.
\texttt{path-mixed}, \texttt{binary-mixed}, \texttt{star-mixed}: fixed
underlying shape (path; balanced binary tree with parent
$\lfloor(v-1)/2\rfloor$; star with hub $0$), only orientations randomized.

\section{Validation methodology}\label{sec:validation}

\paragraph{Exhaustive enumeration oracle.} The CTest suite checks every
algorithm against an oracle that enumerates all closed sets, builds their
affine lines $P-\lambda W$ and computes the exact upper envelope -- a method
structurally unrelated to any algorithm under test. Coverage: every directed
forest with at most four vertices over a finite coefficient grid,
exhaustively; 10\,000 random forests up to 11 vertices; mixed-orientation
path, binary and star trees up to 257 vertices; 6\,000 in-forest and 6\,000
out-forest instances plus 2\,000 larger ones per orientation; and the
structural invariants (the layers partition $V$, each is a closure increment,
thresholds strictly decreasing).

\paragraph{Cross-algorithm differential agreement.} Every campaign runs
several algorithms on the same instance and compares the returned sequences --
partition, thresholds and canonical order, not merely the layer count -- via a
64-bit fingerprint of the canonical serialization. Disagreements are listed in
\texttt{results/processed/mismatches.csv}, never averaged away.
\Cref{tab:correctness} reports the outcome: over 16\,200 instances and
292\,080 timed runs, every pair of algorithms benchmarked on the same instance
returned the identical sequence.

\begin{table}[t]
\centering
\caption{Instances benchmarked and cross-algorithm disagreements per campaign.
Campaign G (\cref{sec:bppf}) is reported separately, \texttt{BPPF} being an
external baseline rather than one of our algorithms.}
\label{tab:correctness}
\small
\tab{correctness}
\end{table}

\paragraph{Sanitizers, CI and exactness.} The suite is built and run on every
push in three configurations: GCC Release, GCC Debug with the address and
undefined-behaviour sanitizers, and Clang Debug. All three pass at the frozen
commit (\path{results/TEST_REPORT_2026-08-31.md}).
\texttt{validate\_instance} rejects any instance whose total absolute profit
or total weight would leave the exact-arithmetic safety margin
($\texttt{INT64\_MAX}/4$), so overflow is refused explicitly rather than
silently wrapped.

\section{Experimental protocol}\label{sec:protocol}

All algorithms are implemented in C++ and run single-threaded on a Linux
machine (Ubuntu 22.04.5, kernel 6.8) with an Intel Core i7-12700 and 32\,GB of
RAM, GCC 11.4.0 at \texttt{-O3}. Every benchmark process is pinned to one
physical core with \texttt{taskset}; the six campaign lanes ran concurrently
on distinct cores, never two processes on one core. Only the algorithm call is
timed: parsing, result hashing and correctness comparison are outside the
timed region. Algorithm order is shuffled independently per repetition.
Repetitions: 11 on the medium random campaign, 3 on the larger and structured
campaigns and on the \texttt{BPPF} comparison, where a single run takes from
tens of milliseconds to tens of seconds. A 300\,s wall-clock timeout and an
8\,GiB memory ceiling (\texttt{ulimit -v}) applied to every run; \emph{no run
of the study hit either}, so no observation is censored.

\paragraph{What we measure, and what ``peak RSS'' means.} Time is wall-clock
time of the algorithm call (\texttt{std::chrono::steady\_clock}); on an
otherwise idle single-threaded run this tracks CPU time closely. Memory is
reported as \emph{peak RSS}, the peak \emph{resident set size} of the process:
the maximum amount of physical RAM the process had mapped at any instant,
excluding pages swapped out or never touched. It is read from
\texttt{getrusage(RUSAGE\_SELF).ru\_maxrss} after each repetition. Two
properties of this counter matter when reading memory figures. First, it is a
\emph{process-lifetime} peak and never decreases, so within one invocation
that benchmarks several algorithms the value attributed to the second
algorithm includes the peak of the first: memory numbers are therefore only
meaningful for single-algorithm invocations, and this report quotes them only
from such runs (\cref{sec:heap}). Second, it counts pages actually resident,
so it includes the allocator's own overhead and is the figure that determines
whether a run survives a memory ceiling -- which is exactly what we want to
report.

\paragraph{Statistics.} Absolute times are medians over repetitions and then
over instances. Comparisons are always \emph{paired}: the ratio is computed
per instance, on the same instance solved by both algorithms, and we report
the median with the interquartile range (IQR) of those per-instance ratios --
never a ratio of aggregate means. The median is used deliberately: it is
invariant under inversion of the comparison, unlike the arithmetic mean of
ratios, which depends on which algorithm sits in the denominator; it is
insensitive to the right tail of the timing noise, which is asymmetric by
construction, since system interference can only slow a run down and never
speed it up; and on every comparison where an advantage is claimed it is the
conservative choice. As a check, on the \texttt{BPPF} comparison at
$n=1\,000$ the three aggregates are $4.5$ (median), $5.8$ (geometric mean)
and $8.7$ (arithmetic mean): the reported figure is the smallest. Where a
population mixes regimes -- as the large random campaign does across four
densities -- no single aggregate describes it, and the breakdowns of
\cref{sec:large} are reported instead of relying on the overall figure.

\paragraph{Known limitation.} The machine has no passwordless root, so the CPU
governor stayed at \texttt{powersave} rather than \texttt{performance}.
Absolute numbers therefore carry ordinary frequency-scaling noise, mitigated
(not eliminated) by median-of-repetitions reporting, core pinning, and by
expressing every headline comparison as a paired same-instance ratio: a
frequency shift during one instance's measurement affects both algorithms
compared on it almost identically.

\section{Campaign design}\label{sec:campaigns}

For every topology and size, 240 instances are generated (four densities, six
coefficient families, ten seeds), or 60 for the structured families, which
have no density parameter. \texttt{mixed-forest}, \texttt{in-forest} and
\texttt{out-forest} instances exist for
$n\in\{100,200,\dots,1\,000\}\cup\{10\,000,20\,000,\dots,100\,000\}$; the
structured families for
$n\in\{100,200,500,1\,000,2\,000,5\,000,10\,000,20\,000,50\,000,100\,000\}$.

\begin{table}[t]
\centering
\caption{The campaigns of the sweep. Two preregistered size cutoffs apply,
both because the asymptotic trend is established below them and running
further would consume hours without adding information:
\texttt{PaC}/\texttt{DPaC} stop at $n=20\,000$ on random forests, and
\texttt{HPaC}/\texttt{DHPaC} stop at $n=20\,000$ on the star class, where
\texttt{PaC} is cheap and is run on the full range instead.}
\label{tab:campaigns}
\begin{adjustbox}{max width=\textwidth}
\small
\begin{tabular}{@{}llll@{}}
\toprule
Campaign & Topology & Sizes $n$ & Algorithms \\
\midrule
A -- correctness & all six & $\le 11$ (exhaustive at 4) & all, vs.\ enumeration oracle \\
B -- medium random & \texttt{mixed-forest} & $100,200,\dots,1\,000$ & \texttt{PaC}, \texttt{DPaC}, \texttt{HPaC}, \texttt{DHPaC}, \texttt{RaC} \\
C -- large random & \texttt{mixed-forest} & $10\,000,\dots,100\,000$ & \texttt{HPaC}, \texttt{DHPaC}, \texttt{RaC}; \texttt{PaC}/\texttt{DPaC} at $10^4,2{\cdot}10^4$ \\
D -- structured & \texttt{path}/\texttt{binary}/\texttt{star-mixed} & $100,\dots,100\,000$ (10 sizes) & \texttt{PaC}, \texttt{DPaC}, \texttt{HPaC}, \texttt{DHPaC}, \texttt{RaC} \\
E -- orientations & \texttt{in-}/\texttt{out-forest} & $100,\dots,100\,000$ (20 sizes) & \texttt{HPaC}, \texttt{DHPaC}, \texttt{HIPaC}/\texttt{HOPaC}, \texttt{RaC} \\
G -- \texttt{BPPF} baseline & \texttt{mixed-forest} & $100,200,\dots,1\,000$ & \texttt{HPaC} vs.\ \texttt{BPPF} \\
\bottomrule
\end{tabular}
\end{adjustbox}
\end{table}

\section{Results}\label{sec:results}

\subsection{Direct scan versus heap on random forests}\label{sec:bcamp}

\Cref{fig:b} reports every algorithm of campaign B on the same instances. The
two families are indistinguishable up to $n\approx 300$ and the heap pays off
from $n\approx 400$ on: the median paired ratio \texttt{PaC}/\texttt{HPaC} is
$0.91$ at $n=100$ -- the direct scan is marginally faster where heap
maintenance does not yet pay for itself -- crosses 1 around $n=400$, and
reaches $2.3$ at $n=1\,000$, $20.3$ at $n=10\,000$ and $42.9$ at
$n=20\,000$. The duals track their primals throughout
(\texttt{DPaC}/\texttt{DHPaC} reaches $54.6$ at $n=20\,000$). The mechanism is
the one anticipated by the manuscript: \texttt{PaC} re-scans every candidate
ratio at every iteration, while \texttt{HPaC} extracts the maximum in
logarithmic time and refreshes only the candidates touched by the last peel or
contraction.

\begin{figure}[t]
\centering
\begin{minipage}[c]{0.47\textwidth}
\centering
\scriptsize
\setlength{\tabcolsep}{3.2pt}
\tab{rep_times_campaign_b}
\end{minipage}%
\hfill
\begin{minipage}[c]{0.49\textwidth}
\centering
\begin{tikzpicture}
\begin{axis}[width=0.97\linewidth, height=5.4cm, xlabel={$n$},
  ylabel={CPU time (ms)}, xmode=log, ymode=log, pcfplot]
\addplot[pac]   table[x=n, y=pac]   {\dat{campaign_b}}; \addlegendentry{\texttt{PaC}}
\addplot[dpac]  table[x=n, y=dpac]  {\dat{campaign_b}}; \addlegendentry{\texttt{DPaC}}
\addplot[hpac]  table[x=n, y=hpac]  {\dat{campaign_b}}; \addlegendentry{\texttt{HPaC}}
\addplot[dhpac] table[x=n, y=dhpac] {\dat{campaign_b}}; \addlegendentry{\texttt{DHPaC}}
\addplot[rac]   table[x=n, y=rac]   {\dat{campaign_b}}; \addlegendentry{\texttt{RaC}}
\end{axis}
\end{tikzpicture}
\end{minipage}
\caption{Median CPU time of all five algorithms on \texttt{mixed-forest}
instances with $n\in\{100,200,\dots,1\,000\}$ (campaign B, 240 instances per
size).}
\label{fig:b}
\end{figure}

The aggregate hides two systematic effects, which the breakdowns make
visible. By density (\cref{tab:b_rho}), all algorithms slow down as the forest
becomes more connected, but not equally: from $\varrho=0.3$ to $\varrho=1.0$
\texttt{HPaC} grows by a factor of about $2$ while \texttt{RaC} grows by less,
which is the medium-size beginning of the density effect that dominates
\cref{sec:large}. By family (\cref{tab:b_family}), the two tie-stressing
families are the fastest for every algorithm, consistently with their smaller
number of closure layers (\cref{tab:layers}); the ordering between algorithms
is the same in every family, so no family reverses a conclusion.

\begin{table}[t]
\centering
\begin{minipage}[t]{0.42\textwidth}
\centering
\caption{Campaign B by density.}
\label{tab:b_rho}
\small
\tab{rep_times_campaign_b_by_rho}
\end{minipage}\hfill
\begin{minipage}[t]{0.53\textwidth}
\centering
\caption{Campaign B by coefficient family.}
\label{tab:b_family}
\scriptsize
\setlength{\tabcolsep}{2.6pt}
\tab{rep_times_campaign_b_by_family}
\end{minipage}
\end{table}

\begin{table}[t]
\centering
\begin{minipage}[t]{0.48\textwidth}
\centering
\caption{Campaign B, paired ratio \texttt{PaC}/\texttt{HPaC} by family:
values above 1 favour \texttt{HPaC}.}
\label{tab:b_ratio_pac}
\scriptsize
\setlength{\tabcolsep}{3.4pt}
\tab{rep_ratio_campaign_b_pac_over_hpac_by_family}
\end{minipage}\hfill
\begin{minipage}[t]{0.48\textwidth}
\centering
\caption{Campaign B, paired ratio \texttt{RaC}/\texttt{HPaC} by family.}
\label{tab:b_ratio_rac}
\scriptsize
\setlength{\tabcolsep}{3.4pt}
\tab{rep_ratio_campaign_b_rac_over_hpac_by_family}
\end{minipage}
\end{table}

\subsection{Heap policy, and why the implementation attains the \texorpdfstring{$\mathcal O(n)$}{O(n)} space bound}
\label{sec:heap}

The candidate heaps admit two natural implementations. The choice is invisible
in the asymptotic time bound and decisive in practice. A push-only
\emph{lazy-deletion} heap re-pushes one entry per incident arc whenever a
vertex's closure sum changes, and discards stale entries only when they
surface at the top: its size is bounded by the number of touch operations,
not by $n$. On a star this is pathological -- every peel or contraction
touches the hub, and each touch re-pushes one entry per incident arc, so the
heap accumulates $\Theta(n^2)$ entries. The implementation used throughout
this study instead rebuilds each heap from the live candidates whenever its
size exceeds a constant multiple of the number of live candidates: an
amortized $\mathcal O(\log n)$ cost per operation that leaves the worst-case
time bound unchanged and makes the memory of the implementation
$\mathcal O(n)$ on \emph{every} input, matching the space bound proved in the
manuscript.

\Cref{tab:heap} quantifies the difference on stars, with single-algorithm,
single-repetition invocations (the only regime in which peak RSS is
attributable to one algorithm, see \cref{sec:protocol}). The lazy heap's peak
resident memory grows by a factor of 100 when $n$ grows by a factor of 10 --
quadratic, as predicted -- and reaches $1.5$\,GiB at $n=10^4$, where the
rebuild policy stays at $8$\,MiB; extrapolating, the lazy policy exhausts an
8\,GiB ceiling by $n\approx 20\,000$. The rebuild policy is also
\emph{faster}, by about $4\times$ here and about $2\times$ on large random
forests, with bit-identical output on every instance tested. This is why the
study reports a single \texttt{HPaC}: the rebuild policy dominates the lazy
one on both axes, and the same policy is used in \texttt{DHPaC}.

\begin{table}[t]
\centering
\caption{The two heap implementations on \texttt{star-mixed} instances
(\texttt{independent-positive}, seed 0, one repetition, Release build). For
reference, \texttt{PaC}, which maintains no candidate structure at all, peaks
at $5.6$\,MiB at $n=10^4$.}
\label{tab:heap}
\small
\begin{tabular}{@{}rrrrr@{}}
\toprule
& \multicolumn{2}{c}{lazy-deletion heap} & \multicolumn{2}{c}{rebuild policy (used here)} \\
\cmidrule(lr){2-3}\cmidrule(l){4-5}
$n$ & time (ms) & peak RSS (MiB) & time (ms) & peak RSS (MiB) \\
\midrule
1\,000 & 139.0 & 15.6 & 64.1 & 3.9 \\
2\,000 & 628.7 & 51.5 & 233.1 & 4.0 \\
5\,000 & 5\,440.9 & 388.0 & 1\,422.1 & 6.4 \\
10\,000 & 27\,020.7 & 1\,541.2 & 6\,396.5 & 8.2 \\
\bottomrule
\end{tabular}
\end{table}

\subsection{Large random forests: size, density and family}\label{sec:large}

On \texttt{mixed-forest} instances up to $n=10^5$ (\cref{fig:c}),
\texttt{HPaC} computes the full parametric solution in about $0.2$\,s and
\texttt{DHPaC} stays within $25\%$ of it at every size. The median paired
ratio \texttt{RaC}/\texttt{HPaC} grows from $4.2$ at $n=10^4$ to $13.3$ at
$n=10^5$, with a widening IQR: \texttt{RaC}'s relative disadvantage on these
instances becomes both larger and more variable as the forests grow.

\begin{figure}[t]
\centering
\begin{minipage}[c]{0.45\textwidth}
\centering
\scriptsize
\setlength{\tabcolsep}{3.2pt}
\tab{rep_times_campaign_c}
\end{minipage}%
\hfill
\begin{minipage}[c]{0.49\textwidth}
\centering
\begin{tikzpicture}
\begin{axis}[width=0.97\linewidth, height=5.4cm, xlabel={$n$},
  ylabel={CPU time (ms)}, xtick={10000,40000,70000,100000},
  minor xtick={20000,30000,50000,60000,80000,90000}, xminorgrids=true,
  minor grid style={gray!15}, scaled x ticks=false,
  /pgf/number format/.cd, fixed, 1000 sep={}, /tikz/.cd, pcfplot]
\addplot[rac]   table[x=n, y=rac]   {\dat{campaign_c}}; \addlegendentry{\texttt{RaC}}
\addplot[dhpac] table[x=n, y=dhpac] {\dat{campaign_c}}; \addlegendentry{\texttt{DHPaC}}
\addplot[hpac]  table[x=n, y=hpac]  {\dat{campaign_c}}; \addlegendentry{\texttt{HPaC}}
\end{axis}
\end{tikzpicture}
\end{minipage}
\caption{Median CPU time on \texttt{mixed-forest} instances with
$n\in\{10\,000,20\,000,\dots,100\,000\}$ (campaign C, 240 instances per size).}
\label{fig:c}
\end{figure}

That aggregate mixes four qualitatively different regimes.
\Cref{tab:c_rho,tab:c_ratio_rho} break it down by density: the median
\texttt{RaC}/\texttt{HPaC} ratio is $12.9$ at $\varrho=0.3$, $17.0$ at
$\varrho=0.6$, $5.0$ at $\varrho=0.9$ and only $2.0$ at $\varrho=1.0$, i.e.\
on single spanning trees, and the IQR shrinks in the same direction
($10.2\to0.4$), so the sparse regime is not only worse for \texttt{RaC} but
far more variable. The intuition is structural: a sparse forest splits into
many small components, on which \texttt{HPaC}'s contractions are almost free,
while \texttt{RaC} still pays its cluster bookkeeping per component; on one
spanning tree the two are closest. Note that the ordering never reverses --
\texttt{HPaC} is faster at every density on these instances -- but the margin
varies by almost an order of magnitude, which is why a single aggregate would
be misleading.

The breakdown by family (\cref{tab:c_family,tab:c_ratio_family}) shows the
opposite picture: it is remarkably flat. Absolute times vary by at most $14\%$
across families for each algorithm, and the paired ratio
\texttt{RaC}/\texttt{HPaC} stays between $5.80$ and $6.11$ in all six. The
coefficient family therefore changes \emph{how much parametric structure}
there is (\cref{tab:layers}) but not \emph{which algorithm wins}: on these
instances the topology, not the coefficients, decides.

\begin{table}[t]
\centering
\begin{minipage}[t]{0.46\textwidth}
\centering
\caption{Campaign C absolute times by density.}
\label{tab:c_rho}
\small
\tab{rep_times_campaign_c_by_rho}
\end{minipage}\hfill
\begin{minipage}[t]{0.50\textwidth}
\centering
\caption{Campaign C, paired ratio \texttt{RaC}/\texttt{HPaC} by density.}
\label{tab:c_ratio_rho}
\small
\tab{campaign_c_rac_over_hpac_by_rho}
\end{minipage}
\end{table}

\begin{table}[t]
\centering
\begin{minipage}[t]{0.50\textwidth}
\centering
\caption{Campaign C absolute times by family.}
\label{tab:c_family}
\scriptsize
\setlength{\tabcolsep}{3.4pt}
\tab{rep_times_campaign_c_by_family}
\end{minipage}\hfill
\begin{minipage}[t]{0.47\textwidth}
\centering
\caption{Campaign C, paired ratio \texttt{RaC}/\texttt{HPaC} by family.}
\label{tab:c_ratio_family}
\scriptsize
\setlength{\tabcolsep}{3.4pt}
\tab{rep_ratio_campaign_c_rac_over_hpac_by_family}
\end{minipage}
\end{table}

\subsection{Structured instance classes}

Paths and balanced binary trees behave like random forests
(\cref{fig:path,fig:binary}): the median ratio \texttt{RaC}/\texttt{HPaC}
lies between $5.8$ and $9.0$ on paths and between $3.2$ and $3.9$ on binary
trees, from $n=100$ to $n=100\,000$, with no trend towards a crossover. On
both classes \texttt{PaC} behaves as its $\mathcal O(n^2)$ bound predicts and
is left far behind by $n=10^4$.

\begin{figure}[t]
\centering
\begin{minipage}[c]{0.45\textwidth}
\centering
\scriptsize
\setlength{\tabcolsep}{3.2pt}
\tab{rep_times_campaign_d_path}
\end{minipage}%
\hfill
\begin{minipage}[c]{0.49\textwidth}
\centering
\begin{tikzpicture}
\begin{axis}[width=0.97\linewidth, height=5.4cm, xlabel={$n$},
  ylabel={CPU time (ms)}, xmode=log, ymode=log, pcfplot]
\addplot[pac]  table[x=n, y=pac]  {\dat{campaign_d_path}}; \addlegendentry{\texttt{PaC}}
\addplot[hpac] table[x=n, y=hpac] {\dat{campaign_d_path}}; \addlegendentry{\texttt{HPaC}}
\addplot[rac]  table[x=n, y=rac]  {\dat{campaign_d_path}}; \addlegendentry{\texttt{RaC}}
\end{axis}
\end{tikzpicture}
\end{minipage}
\caption{Median CPU time on \texttt{path-mixed} instances (campaign D, 60
instances per size).}
\label{fig:path}
\end{figure}

\begin{figure}[t]
\centering
\begin{minipage}[c]{0.45\textwidth}
\centering
\scriptsize
\setlength{\tabcolsep}{3.2pt}
\tab{rep_times_campaign_d_binary}
\end{minipage}%
\hfill
\begin{minipage}[c]{0.49\textwidth}
\centering
\begin{tikzpicture}
\begin{axis}[width=0.97\linewidth, height=5.4cm, xlabel={$n$},
  ylabel={CPU time (ms)}, xmode=log, ymode=log, pcfplot]
\addplot[pac]  table[x=n, y=pac]  {\dat{campaign_d_binary}}; \addlegendentry{\texttt{PaC}}
\addplot[hpac] table[x=n, y=hpac] {\dat{campaign_d_binary}}; \addlegendentry{\texttt{HPaC}}
\addplot[rac]  table[x=n, y=rac]  {\dat{campaign_d_binary}}; \addlegendentry{\texttt{RaC}}
\end{axis}
\end{tikzpicture}
\end{minipage}
\caption{Median CPU time on \texttt{binary-mixed} instances (campaign D, 60
instances per size).}
\label{fig:binary}
\end{figure}

Stars invert the picture, increasingly sharply with $n$ (\cref{fig:star}):
\texttt{RaC} is already faster than \texttt{HPaC} at $n=100$ (median ratio
$0.75$) and about $250$ times faster at $n=20\,000$. The mechanism is the one
of \cref{sec:heap} seen from the algorithmic side: every peel or contraction
touches the hub and invalidates the candidate ratios of a constant fraction of
the remaining leaves, so \emph{any} heap-based schedule pays $\Theta(n)$
maintenance per event, while \texttt{RaC}'s rake operations absorb all leaves
in $\mathcal O(\log n)$ rounds regardless of the hub's update history. The
direct scan is far less affected, since it maintains no candidate structure:
\texttt{PaC} stays about one order of magnitude below \texttt{HPaC} over the
whole range and is overtaken by \texttt{RaC} only between $n=1\,000$ and
$n=2\,000$. This is the one class of the study where the better worst-case
bound of \texttt{RaC} is visible in practice, and there it is unambiguous.
The effect is uniform across coefficient families
(\cref{tab:star_ratio_family}): the median \texttt{RaC}/\texttt{HPaC} ratio lies
between $0.060$ and $0.064$ in all six, so it is the topology alone that
produces it.

\begin{figure}[t]
\centering
\begin{minipage}[c]{0.45\textwidth}
\centering
\scriptsize
\setlength{\tabcolsep}{3.2pt}
\tab{rep_times_campaign_d_star}
\end{minipage}%
\hfill
\begin{minipage}[c]{0.49\textwidth}
\centering
\begin{tikzpicture}
\begin{axis}[width=0.97\linewidth, height=5.4cm, xlabel={$n$},
  ylabel={CPU time (ms)}, xmode=log, ymode=log, pcfplot]
\addplot[hpac] table[x=n, y=hpac] {\dat{campaign_d_star}}; \addlegendentry{\texttt{HPaC}}
\addplot[pac]  table[x=n, y=pac]  {\dat{campaign_d_star}}; \addlegendentry{\texttt{PaC}}
\addplot[rac]  table[x=n, y=rac]  {\dat{campaign_d_star}}; \addlegendentry{\texttt{RaC}}
\end{axis}
\end{tikzpicture}
\end{minipage}
\caption{Median CPU time on \texttt{star-mixed} instances (campaign D, 60
instances per size). The heap-based algorithms stop at the preregistered
cutoff $n=20\,000$; \texttt{PaC} and \texttt{RaC} run the full range.}
\label{fig:star}
\end{figure}

\begin{table}[t]
\centering
\begin{minipage}[t]{0.48\textwidth}
\centering
\caption{Stars, paired ratio \texttt{RaC}/\texttt{HPaC} by family: values
below 1 favour \texttt{RaC}.}
\label{tab:star_ratio_family}
\scriptsize
\setlength{\tabcolsep}{3.4pt}
\tab{rep_ratio_campaign_d_star_rac_over_hpac_by_family}
\end{minipage}\hfill
\begin{minipage}[t]{0.48\textwidth}
\centering
\caption{Stars, paired ratio \texttt{RaC}/\texttt{PaC} by size: the crossover
between the two non-heap-limited algorithms.}
\label{tab:star_ratio_pac}
\scriptsize
\tab{campaign_d_star_rac_over_pac}
\end{minipage}
\end{table}

\subsection{Specialized variants for single orientations}

On the orientation each is built for, the specialized $\mathcal O(n\log n)$
variants take a stable fraction of the general algorithm's time
(\cref{fig:ein,fig:eout}): the median ratio \texttt{HIPaC}/\texttt{HPaC} lies
between $0.52$ and $0.75$ on in-forests and \texttt{HOPaC}/\texttt{HPaC}
between $0.57$ and $0.78$ on out-forests, with no erosion up to $n=10^5$ and
no family in which the advantage disappears
(\cref{tab:e_ratio_family_in,tab:e_ratio_family_out}). The gain reflects the
structural reason given in the manuscript: on a single orientation the minimal
preceding (respectively succeeding) set of every arc collapses to a singleton,
so one heap value per vertex suffices and no closure-sum propagation is needed.
On the same instances \texttt{RaC} loses ground as $n$ grows, from about
$2.8\times$ \texttt{HPaC} at $n\le 1\,000$ to $12.8\times$ at $n=10^5$,
mirroring the trend on general mixed forests.

\begin{figure}[t]
\centering
\begin{minipage}[c]{0.45\textwidth}
\centering
\scriptsize
\setlength{\tabcolsep}{3.2pt}
\tab{rep_times_campaign_e_in}
\end{minipage}%
\hfill
\begin{minipage}[c]{0.49\textwidth}
\centering
\begin{tikzpicture}
\begin{axis}[width=0.97\linewidth, height=5.4cm, xlabel={$n$},
  ylabel={CPU time (ms)}, xmode=log, ymode=log, pcfplot]
\addplot[hpac]  table[x=n, y=hpac]  {\dat{campaign_e_in}}; \addlegendentry{\texttt{HPaC}}
\addplot[hipac] table[x=n, y=hipac] {\dat{campaign_e_in}}; \addlegendentry{\texttt{HIPaC}}
\addplot[rac]   table[x=n, y=rac]   {\dat{campaign_e_in}}; \addlegendentry{\texttt{RaC}}
\end{axis}
\end{tikzpicture}
\end{minipage}
\caption{Median CPU time on \texttt{in-forest} instances (campaign E, 240
instances per size).}
\label{fig:ein}
\end{figure}

\begin{figure}[t]
\centering
\begin{minipage}[c]{0.45\textwidth}
\centering
\scriptsize
\setlength{\tabcolsep}{3.2pt}
\tab{rep_times_campaign_e_out}
\end{minipage}%
\hfill
\begin{minipage}[c]{0.49\textwidth}
\centering
\begin{tikzpicture}
\begin{axis}[width=0.97\linewidth, height=5.4cm, xlabel={$n$},
  ylabel={CPU time (ms)}, xmode=log, ymode=log, pcfplot]
\addplot[hpac]  table[x=n, y=hpac]  {\dat{campaign_e_out}}; \addlegendentry{\texttt{HPaC}}
\addplot[hopac] table[x=n, y=hopac] {\dat{campaign_e_out}}; \addlegendentry{\texttt{HOPaC}}
\addplot[rac]   table[x=n, y=rac]   {\dat{campaign_e_out}}; \addlegendentry{\texttt{RaC}}
\end{axis}
\end{tikzpicture}
\end{minipage}
\caption{Median CPU time on \texttt{out-forest} instances (campaign E, 240
instances per size).}
\label{fig:eout}
\end{figure}

\begin{table}[t]
\centering
\begin{minipage}[t]{0.48\textwidth}
\centering
\caption{In-forests, \texttt{HIPaC}/\texttt{HPaC} by family.}
\label{tab:e_ratio_family_in}
\scriptsize
\setlength{\tabcolsep}{3.4pt}
\tab{rep_ratio_campaign_e_in_hipac_over_hpac_by_family}
\end{minipage}\hfill
\begin{minipage}[t]{0.48\textwidth}
\centering
\caption{Out-forests, \texttt{HOPaC}/\texttt{HPaC} by family.}
\label{tab:e_ratio_family_out}
\scriptsize
\setlength{\tabcolsep}{3.4pt}
\tab{rep_ratio_campaign_e_out_hopac_over_hpac_by_family}
\end{minipage}
\end{table}

We record explicitly that this advantage is smaller than the one reported in
the first version of the manuscript, where the specialized variants were
$4$--$10$ times faster than the general algorithm. The direction of the change
is expected -- the general \texttt{HPaC} is substantially faster than it was,
for the reasons of \cref{sec:heap} -- but its size is only partly accounted
for by the heap policy: a direct comparison of the two heap implementations on
these instances shows a factor of $1.3$--$1.4$, not $3$--$4$. The remainder is
attributable to the code being an independent rewrite rather than a
refactoring, and to instances regenerated under different coefficient
families. We therefore treat absolute times across the two versions as
incomparable, and only trends as comparable.

\subsection{Comparison with parametric pseudoflow}\label{sec:bppf}

\texttt{BPPF} is the unmodified upstream implementation of the parametric
pseudoflow method, applicable to arbitrary precedence graphs. It evaluates
minimum cuts at user-supplied parameter values, so a comparison requires
deciding what to supply, and we drive it in the setting most favourable to it,
identical to the methodology of the manuscript's first version: for each
instance, one single \texttt{BPPF} process receives, in \texttt{BPPF}'s native
affine (two-numbers-per-arc) capacity format, the $k+1$ parameter values that
bracket all $k$ breakpoints. \texttt{BPPF} is thus spared the search for the
breakpoints, and one call certifies the entire parametric solution. Its own
internal solve timer (input parsing excluded) is compared against
\texttt{HPaC}'s in-process time for the full sweep on the same instance; the
fixed-point precision is $10^{-6}$.

On all 2\,400 instances with $n\le 1\,000$, the median paired ratio
\texttt{BPPF}/\texttt{HPaC} grows from $1.5$ at $n=100$ to $4.5$ at
$n=1\,000$, with an overall median of $3.3$. Two qualifications belong with
this number. First, it is measured on forests with $n\le 1\,000$, and
\texttt{BPPF} solves a strictly more general problem, so it says nothing about
\texttt{BPPF} outside this class. Second, the closures returned by the two
methods agree at every probe on every instance, with one exception class: on
16 of the 2\,400 instances \texttt{BPPF} merges two consecutive closure layers
whose thresholds differ by less than its fixed-point precision, and therefore
does not recover the exact optimal layer sequence. All 16 belong to
\texttt{near-ties} and \texttt{exact-ties}, the families designed precisely to
place thresholds at that distance. The algorithms of this study are
unaffected, since they compare integers exactly; this is the practical
consequence of exact arithmetic, and the reason the comparison reports
agreement separately from timing.

An earlier, discarded variant of this comparison drove \texttt{BPPF} once per
breakpoint, one process spawn per call. That measurement is reported nowhere:
the resulting ratio is dominated by process-creation overhead and says nothing
about either algorithm.

\section{Computational conclusions}

Correctness is uniform across the study: 16\,200 instances, 292\,080 timed
runs, seven algorithms, zero cross-algorithm disagreements, on top of the
exhaustive enumeration oracle on small instances and the sanitizer builds. No
run hit the time or memory limit, so no result is censored or imputed.

On the instances tested, \texttt{HPaC} is the algorithm of choice on
mixed-orientation random forests, paths and balanced binary trees, faster than
\texttt{RaC} by a factor between $2$ and $13$ depending on size and density,
and faster than the direct scan from $n\approx 400$ on. \texttt{RaC} is the
algorithm of choice on stars, where its advantage grows with $n$ and reaches
about $250\times$ at $n=20\,000$ -- the one class where its better worst-case
bound is visible in practice, for a structural reason that also explains why
the direct scan resists there better than any heap-based schedule.
\texttt{HIPaC} and \texttt{HOPaC} deliver a stable $1.3$--$1.9\times$ speedup
over the general algorithm when the orientation is known in advance. Against a
general parametric pseudoflow implementation on forests with $n\le 1\,000$,
\texttt{HPaC} is faster by a median factor of $3.3$ even when the pseudoflow
solver is handed the breakpoint locations, and unlike it recovers the exact
layer sequence on every instance.

Two structural facts organize all of the above. The topology decides which
algorithm wins -- density moves the margin by an order of magnitude and the
star class inverts the ordering entirely -- while the coefficient family
decides only how much parametric structure exists, moving absolute times but
never the ranking. And the choice of data structure inside one algorithm can
matter as much as the choice of algorithm: the two heap policies of
\cref{sec:heap} differ by $4\times$ in time and two orders of magnitude in
memory, at identical output and identical asymptotic bounds.

None of these statements is a claim of general dominance: each is scoped to
the topologies, sizes and coefficient families in which it was measured, and
the star class is the standing reminder that the ordering can invert
completely when the topology changes.

\section*{Code and data availability}

The code, instance generators, raw and processed benchmark data and the
analysis pipeline are available at
\url{https://github.com/fabiofurini/parametric-closure-forests}, release
\texttt{v0.3.0}. The release attaches the full instance archive (split under
GitHub's asset size limit) and the raw and processed benchmark data, with
SHA-256 checksums. Every table and every plot in this report is generated from
\texttt{results/raw/*.csv} by \texttt{tools/build\_report.sh} and
\texttt{tools/emit\_report\_tables.py}; regenerating them from the same raw
data reproduces this report exactly. The pseudoflow baseline is the unmodified
upstream implementation available at
\url{https://github.com/hochbaumGroup/Bounded-precision-simple-parametric.git}.

\end{document}
